\FloatBarrier
\section{Utility tolerance, case floors, and exchange uncertainty}
\label{app:v8}
These additional analyses are retrospective and use the existing paired fits, frozen candidate scores and previously evaluated periods. Their specifications, input hashes and code hashes were recorded before the new computations. The scientific choices were informed by earlier results and review; this chronology is not prospective preregistration. No new models or external data were introduced. All grid settings and no-choice conditions are retained in the accompanying source, under \texttt{evidence/revision\_v8\_*}.

\subsection{Conditional uncertainty of the observed exchanges}
\label{app:v8-exchange}
We keep every original relative-.95 policy choice, cap and threshold fixed. For each arm/head cell, exposure-choice minus case-choice differences define $Q=\Delta d/(-\Delta c)$. We draw 4,000 item-cluster multiplicity vectors at seed 20260915, shared across both arms and all head cells. Each vector retains the item's stores and dates. All case-loss denominators remain strictly positive in every draw; the smallest across the eight cells is 13.5332 points. Thus ordinary percentile summaries of these fixed-policy ratios are defined. Table~\ref{tab:v8-exchange} gives every point estimate and central 95\% interval. The common multiplicities also provide complete-minus-censored differences, with pointwise and four-contrast-adjusted intervals in Table~\ref{tab:v8-exchange-difference}.
\begin{table}[t]\centering\small
\caption{All eight original policy exchanges $\Delta d/(-\Delta c)$ at each arm\textquotesingle s primary relative cap. Central 95\% percentile intervals use 4,000 item resamples and condition on fits, caps, thresholds and choices. All case-loss denominators are positive in every draw; these intervals omit calibration reselection.}\label{tab:v8-exchange}
\begin{tabular}{llrl}\toprule
Arm & Heads $e/w$ & Exchange & Conditional interval\\\midrule
Censored & 220/220 & 0.1764 & [0.1155, 0.2472] \\
Censored & 660/220 & 0.2040 & [0.1374, 0.2811] \\
Censored & 220/660 & 0.2130 & [0.1508, 0.2835] \\
Censored & 660/660 & 0.2034 & [0.1393, 0.2781] \\
Complete & 220/220 & 0.6150 & [0.4829, 0.7779] \\
Complete & 660/220 & 0.3316 & [0.2548, 0.4302] \\
Complete & 220/660 & 0.2835 & [0.2181, 0.3621] \\
Complete & 660/660 & 0.3061 & [0.2316, 0.4019] \\
\bottomrule\end{tabular}\end{table}

\begin{table}[t]\centering\small
\caption{Complete-minus-censored exchange differences, with shared item multiplicities across arms. The four-contrast intervals use empirical .00625/.99375 quantiles. Adjustment gives approximate conditional summaries, not simultaneous population certification or a causal effect of censoring.}\label{tab:v8-exchange-difference}
\begin{tabular}{lrll}\toprule
Heads $e/w$ & Difference & Pointwise interval & Four-contrast interval\\\midrule
220/220 & 0.4387 & [0.3577, 0.5403] & [0.3398, 0.5750] \\
660/220 & 0.1277 & [0.0912, 0.1658] & [0.0841, 0.1762] \\
220/660 & 0.0705 & [0.0459, 0.0951] & [0.0400, 0.1015] \\
660/660 & 0.1027 & [0.0709, 0.1377] & [0.0626, 0.1479] \\
\bottomrule\end{tabular}\end{table}

The complete-sales point range is .283511--.615044; the censored range is .176389--.212997. All four matched-head point differences and adjusted conditional intervals are positive. This characterizes the saved exchanges at their different arm-relative caps, not a general or causal efficiency claim. At the complete 220/220 reference, the originally selected Ratio branch has $Q=.615044$ and interval [.482931,.777873]. Replacing it with the existing Descending candidate, keeping the Excess case policy, gives -35.1632 case/+10.0836 exposure points and $Q=.286766$ [.218574,.368456]. This second calculation is a fixed counterfactual candidate comparison, not a new selection experiment.

These intervals omit calibration reselection, model fitting and cap estimation uncertainty, and item resampling does not remove shared store/calendar shocks. In particular, a 48.5\% family frequency in the earlier calibration diagnostic is not a literal probability that the observed reference was generated by a coin toss. We have not estimated a selection-inclusive exchange distribution and make no claim that it is bimodal. Fixed-policy uncertainty is estimable; it must not be substituted for uncertainty in the full pipeline.

\subsection{Margins at stricter relative caps}
\label{app:v8-strict}
We independently reconstruct all 30 original calibration candidates at relative factors .85/.90/.95, for the 220/220 reference in both arms. At .85 and .90, Descending has no feasible threshold in either arm. Its Ratio comparison is therefore undefined, not a numerical margin against zero coverage. In complete sales, Ratio is selected and exceeds the strongest feasible other candidate, Mixed, by 4.0141 and 3.4327 minimum-calibration-exposure points. In censored sales, .85 selects Mixed with Ratio infeasible; .90 selects Ratio, ahead of Mixed by 2.5829 points (Table~\ref{tab:v8-strict}). Complete .85 Ratio has minimum case coverage 35.0362\%, very close to the .35 floor. These are menu-feasibility and calibration-utility observations. They do not isolate the causal value of the error head, establish population-optimal ranking, or establish robustness to a stricter floor.
\begin{table}[t]\centering\scriptsize
\caption{Original 220/220 calibration menus at three relative caps. Margins are minimum-calibration-exposure differences in pp. R/D denote Ratio/Descending. Undefined comparisons are shown as -- when a candidate is infeasible; zero accepted mass is not a feasible competing policy. The strongest feasible other candidate is Mixed at .85/.90 and Descending at .95.}\label{tab:v8-strict}
\begin{tabular}{lrrlllrr}\toprule
Arm & Cap factor & Eligible & Winner & R feasible & D feasible & R--D & R--Other\\\midrule
Censored & 0.85 & 2 & Mixed & no & no & -- & -- \\
Censored & 0.90 & 3 & Ratio & yes & no & -- & 2.5829 \\
Censored & 0.95 & 4 & Descending & yes & yes & -0.5442 & -0.5442 \\
Complete & 0.85 & 3 & Ratio & yes & no & -- & 4.0141 \\
Complete & 0.90 & 3 & Ratio & yes & no & -- & 3.4327 \\
Complete & 0.95 & 4 & Ratio & yes & yes & 0.0270 & 0.0270 \\
\bottomrule\end{tabular}\end{table}


\FloatBarrier
\subsection{A declared exposure tolerance and secondary case utility}
\label{app:v8-tolerance}
The specification fixes the illustrative primary allowance at $\varepsilon=.005$ in coverage fractions (.5 percentage points), alongside the full grid $\{0,.001,.0025,.005,.01\}$. This is a disclosed choice informed by the earlier audit, not a data-independent operational standard. It allows at most the stated loss in the primary calibration utility in exchange for case coverage. We neither select the allowance from evaluation results nor interpret it as a statistical confidence threshold.

For every one of the eight original paired cells at relative .95, original candidate thresholds and feasibility remain unchanged. Equation~\eqref{eq:tolerance} maximizes minimum A/B case coverage among candidates whose minimum A/B exposure coverage is within $\varepsilon$ of the best, then maximizes exposure for any exact case tie, then uses the fixed original method order. An empty eligible menu returns no choice. Case utility remains the original exact maximizer. Since the original exposure winner belongs to the band, the new minimum-calibration case coverage cannot be smaller, and the calibration exposure concession is at most $\varepsilon$. These are properties of the selection rule on the finite calibration menu, not generalization bounds or proof of better population utility.

All 40 cell/allowance choices were frozen before rereading evaluation outcomes. Table~\ref{tab:v8-tolerance-grid} gives the whole selection grid; Table~\ref{tab:v8-tolerance} gives all eight cells at .5 points, including unchanged censored results. All complete-sales cells then select Ratio. The three changed cells sacrifice .0938--.3884 calibration exposure points, while gaining 13.62--24.11 evaluation case points and losing .2506--.6208 evaluation exposure points relative to their original exposure choice. Complete 660/220 exceeds the .5 allowance on evaluation. The procedure preserves the original candidate constraints but has no additional evaluation-risk certification; it is separate from the A-design/B-screen experiment.
\begin{table}[t]\centering\scriptsize
\caption{Complete tolerance grid: exposure-rule choices with original thresholds held fixed. Column values are exposure-coverage allowances in pp. Case choices remain Excess throughout. This is an explicit utility preference, not a confidence region. All candidate metrics and 40 condition outcomes are retained in the accompanying results.}\label{tab:v8-tolerance-grid}
\begin{tabular}{lllllll}\toprule
Arm & Heads & 0 & .1 & .25 & .5 & 1.0\\\midrule
Censored & 220/220 & Descending & Descending & Descending & Descending & Ratio \\
Censored & 660/220 & Descending & Descending & Descending & Descending & Descending \\
Censored & 220/660 & Descending & Descending & Descending & Descending & Ratio \\
Censored & 660/660 & Descending & Descending & Descending & Descending & Descending \\
Complete & 220/220 & Ratio & Ratio & Ratio & Ratio & Ratio \\
Complete & 660/220 & Descending & Descending & Descending & Ratio & Ratio \\
Complete & 220/660 & Descending & Ratio & Ratio & Ratio & Ratio \\
Complete & 660/660 & Descending & Descending & Descending & Ratio & Ratio \\
\bottomrule\end{tabular}\end{table}

\begin{table}[t]\centering\scriptsize
\caption{All eight cells at the declared .5 pp exposure tolerance. Case choice remains Excess. Cal. cost is the minimum-calibration-exposure concession; Later gain/cost compare the tolerance-selected exposure policy with the original exposure choice. Final columns compare the new exposure choice with the unchanged case choice. The .621 pp Later cost in complete 660/220 exceeds the calibration allowance.}\label{tab:v8-tolerance}
\begin{tabular}{llllrrrr}\toprule
Arm & Heads & Original/new & Cal. cost & Later $c$ gain & Later $d$ cost & $\Delta c$ & $\Delta d$\\\midrule
Censored & 220/220 & Descending/Descending & 0.000 & 0.00 & -0.000 & -39.95 & 7.05 \\
Censored & 660/220 & Descending/Descending & 0.000 & 0.00 & -0.000 & -40.66 & 8.29 \\
Censored & 220/660 & Descending/Descending & 0.000 & 0.00 & -0.000 & -37.94 & 8.08 \\
Censored & 660/660 & Descending/Descending & 0.000 & 0.00 & -0.000 & -40.04 & 8.14 \\
Complete & 220/220 & Ratio/Ratio & 0.000 & 0.00 & -0.000 & -15.81 & 9.72 \\
Complete & 660/220 & Descending/Ratio & 0.388 & 24.11 & 0.621 & -10.41 & 10.83 \\
Complete & 220/660 & Descending/Ratio & 0.094 & 13.62 & 0.251 & -21.23 & 9.63 \\
Complete & 660/660 & Descending/Ratio & 0.360 & 19.00 & 0.476 & -15.98 & 10.23 \\
\bottomrule\end{tabular}\end{table}


For stability, we apply every allowance to the \emph{same} earlier 200 threshold-redesign draws for complete 220/220 and Bike. Each draw already recalibrates all ranking thresholds. Only its final menu choice is reanalyzed, conditional on fixed fits, cap, $T$ and draw identities; no new resampling or evaluation access enters this comparison. Exact zero tolerance reproduces the original frequencies. At .5 points, complete 220/220 selects Ratio in 194/200, Mixed in 3/200 and Descending in 3/200; Bike selects Ratio in 172/200 and Mixed in 28/200. The complete setting's Ratio frequency then falls to 167/200 at 1 point, with Mixed in 32/200 and Descending in 1/200 (Table~\ref{tab:v8-tolerance-frequency}). Thus the band changes the trade-off and need not monotonically stabilize a named rule. In the complete cells it removes the \emph{point-estimate family switch} across head counts; coverage values still depend on the fitted heads, and only the reference cell has this redesign-resampling comparison.
\begin{table}[t]\centering\scriptsize
\caption{Exposure-rule frequency (\%) after applying each tolerance to the same 200 threshold-redesign draws from the earlier audit. Fitted scores, cap, $T$, and draw identities are unchanged. The .5 pp rule makes Ratio more frequent in the complete-sales reference and less frequent in Bike; a wider band need not monotonically favor Ratio.}\label{tab:v8-tolerance-frequency}
\begin{tabular}{llrrrrrr}\toprule
Setting & Tolerance pp & Error & Excess & Mixed & Ratio & Desc. & None\\\midrule
Complete 220/220 & 0.00 & 0.0 & 0.0 & 0.0 & 51.5 & 48.5 & 0.0 \\
Complete 220/220 & 0.10 & 0.0 & 0.0 & 0.0 & 69.0 & 31.0 & 0.0 \\
Complete 220/220 & 0.25 & 0.0 & 0.0 & 0.5 & 83.5 & 16.0 & 0.0 \\
Complete 220/220 & 0.50 & 0.0 & 0.0 & 1.5 & 97.0 & 1.5 & 0.0 \\
Complete 220/220 & 1.00 & 0.0 & 0.0 & 16.0 & 83.5 & 0.5 & 0.0 \\
Bike & 0.00 & 0.0 & 0.0 & 4.0 & 96.0 & 0.0 & 0.0 \\
Bike & 0.10 & 0.0 & 0.0 & 4.5 & 95.5 & 0.0 & 0.0 \\
Bike & 0.25 & 0.0 & 0.0 & 8.0 & 92.0 & 0.0 & 0.0 \\
Bike & 0.50 & 0.0 & 0.0 & 14.0 & 86.0 & 0.0 & 0.0 \\
Bike & 1.00 & 0.0 & 0.0 & 34.5 & 65.5 & 0.0 & 0.0 \\
\bottomrule\end{tabular}\end{table}

The source retains every candidate's calibration and evaluation coverage and risk, every cell/allowance choice, all signed changes, and all 2,000 draw/allowance choices. Independent lexicographic reconstruction, exact zero-tolerance replay, concession inequalities, whole-grid checks, and boundary toy examples validate the implementation. Exact ties use a numerical tolerance of $10^{-12}$ only for floating-point comparisons; .5 pp is the substantive allowance.

\FloatBarrier
\subsection{Separate-window screening across case floors}
\label{app:v8-floor}
The protocol specifies floors .35/.40/.45, all four complete-sales head pairs and all five ranking families: 60 candidate conditions in total. Reporting cap .6326834234, score cap $r$, design cap $.95r=.6010492522$, $T$, windows and fitted arrays are unchanged. Each condition searches the largest feasible whole-score-tie threshold on A, then is screened on B using 4,000 item resamples at seed 20260912. The primary sensitivity screen uses tail .05/60 across the entire grid; the inherited .05/20 screen is also reported to reproduce the original design. Utility uses A alone among B passers, with the original exact criterion and tie order. The exposure-tolerance rule is not applied in this floor experiment. All choices were frozen before the existing Later arrays were reopened.

For a fixed ranking, accepted case coverage increases with its threshold. Raising the floor therefore retains the previous largest risk-feasible threshold if it clears the new floor, or makes the family infeasible; a lower threshold cannot restore the larger case floor. Independent tie-endpoint enumeration verifies this behavior. Descending's maximum A case coverage is 35.0483\% (exposure head 220) or 35.0938\% (660). It already fails a .36 floor and is infeasible at both .40 and .45. Ratio remains feasible throughout the declared grid, but at 220/660 its A coverage is only 45.0574\%. Other larger floors need not preserve the observed choices.
\begin{table}[t]\centering\scriptsize
\caption{All 60 floor/head/ranking conditions: B signed-excess upper summaries at tail .05/60, per item. -- means no A-feasible threshold. All defined entries are negative and pass B, so B eliminates no additional candidate. The floor change affects candidate eligibility on A. Full threshold and mass records are retained in the accompanying source.}\label{tab:v8-floor-screen}
\begin{tabular}{llrrrrr}\toprule
Floor & Heads & Error & Excess & Mixed & Ratio & Descending\\\midrule
0.35 & 220/220 & -- & -1.630 & -2.907 & -3.309 & -3.130 \\
0.40 & 220/220 & -- & -1.630 & -2.907 & -3.309 & -- \\
0.45 & 220/220 & -- & -1.630 & -2.907 & -3.309 & -- \\
0.35 & 660/220 & -- & -1.718 & -2.664 & -3.037 & -3.130 \\
0.40 & 660/220 & -- & -1.718 & -2.664 & -3.037 & -- \\
0.45 & 660/220 & -- & -1.718 & -2.664 & -3.037 & -- \\
0.35 & 220/660 & -- & -1.657 & -2.843 & -3.228 & -3.192 \\
0.40 & 220/660 & -- & -1.657 & -2.843 & -3.228 & -- \\
0.45 & 220/660 & -- & -1.657 & -2.843 & -3.228 & -- \\
0.35 & 660/660 & -- & -1.485 & -2.594 & -3.311 & -3.192 \\
0.40 & 660/660 & -- & -1.485 & -2.594 & -3.311 & -- \\
0.45 & 660/660 & -- & -1.485 & -2.594 & -3.311 & -- \\
\bottomrule\end{tabular}\end{table}

Across all grid conditions, 40 candidates are nonempty: 16 at .35 and 12 each at .40/.45. Every one passes B under both family adjustments. The changed choice therefore results from the A feasibility restriction, not from additional B rejection. The .35/.05-over-20 thresholds, B checks, selected policies, Later point metrics and risk intervals reproduce the original screen. At .40/.45, all four cells select Excess for cases and Ratio for exposure. Table~\ref{tab:v8-floor-later} gives the full 24 selected records under the 60-condition family, including duplicates; all 48 records counting both family adjustments are saved in the source.
\begin{table}[t]\centering\scriptsize
\caption{All 24 selected condition records under the 60-candidate screen family on the previously used Later window. $c,d$ are percentages; changes are pp relative to the same head/utility at floor .35. Rows include duplicated policies and are not independent replications. $U_{60}$ is an approximate one-sided .9991667 quantile based on 4,000 draws.}\label{tab:v8-floor-later}
\begin{tabular}{llllrrrrrr}\toprule
Floor & Heads & Utility & Rule & $c$ & $d$ & Risk & $U_{60}$ & Change $c$ & Change $d$\\\midrule
0.35 & 220/220 & c & Excess & 73.55 & 70.81 & 0.59428 & 0.62677 & 0.00 & 0.00 \\
0.35 & 220/220 & d & Descending & 37.16 & 83.89 & 0.59787 & 0.62793 & 0.00 & 0.00 \\
0.40 & 220/220 & c & Excess & 73.55 & 70.81 & 0.59428 & 0.62677 & 0.00 & 0.00 \\
0.40 & 220/220 & d & Ratio & 50.36 & 82.62 & 0.59658 & 0.62613 & 13.20 & -1.26 \\
0.45 & 220/220 & c & Excess & 73.55 & 70.81 & 0.59428 & 0.62677 & 0.00 & 0.00 \\
0.45 & 220/220 & d & Ratio & 50.36 & 82.62 & 0.59658 & 0.62613 & 13.20 & -1.26 \\
0.35 & 660/220 & c & Excess & 72.99 & 69.60 & 0.59446 & 0.62683 & 0.00 & 0.00 \\
0.35 & 660/220 & d & Descending & 37.16 & 83.89 & 0.59787 & 0.62793 & 0.00 & 0.00 \\
0.40 & 660/220 & c & Excess & 72.99 & 69.60 & 0.59446 & 0.62683 & 0.00 & 0.00 \\
0.40 & 660/220 & d & Ratio & 56.18 & 82.32 & 0.59725 & 0.62691 & 19.02 & -1.57 \\
0.45 & 660/220 & c & Excess & 72.99 & 69.60 & 0.59446 & 0.62683 & 0.00 & 0.00 \\
0.45 & 660/220 & d & Ratio & 56.18 & 82.32 & 0.59725 & 0.62691 & 19.02 & -1.57 \\
0.35 & 220/660 & c & Excess & 73.27 & 71.19 & 0.59484 & 0.62706 & 0.00 & 0.00 \\
0.35 & 220/660 & d & Descending & 37.25 & 83.96 & 0.59794 & 0.62802 & 0.00 & 0.00 \\
0.40 & 220/660 & c & Excess & 73.27 & 71.19 & 0.59484 & 0.62706 & 0.00 & 0.00 \\
0.40 & 220/660 & d & Ratio & 46.49 & 82.78 & 0.59704 & 0.62689 & 9.24 & -1.17 \\
0.45 & 220/660 & c & Excess & 73.27 & 71.19 & 0.59484 & 0.62706 & 0.00 & 0.00 \\
0.45 & 220/660 & d & Ratio & 46.49 & 82.78 & 0.59704 & 0.62689 & 9.24 & -1.17 \\
0.35 & 660/660 & c & Excess & 73.55 & 70.43 & 0.59550 & 0.62756 & 0.00 & 0.00 \\
0.35 & 660/660 & d & Descending & 37.25 & 83.96 & 0.59794 & 0.62802 & 0.00 & 0.00 \\
0.40 & 660/660 & c & Excess & 73.55 & 70.43 & 0.59550 & 0.62756 & 0.00 & 0.00 \\
0.40 & 660/660 & d & Ratio & 51.20 & 82.54 & 0.59706 & 0.62672 & 13.95 & -1.42 \\
0.45 & 660/660 & c & Excess & 73.55 & 70.43 & 0.59550 & 0.62756 & 0.00 & 0.00 \\
0.45 & 660/660 & d & Ratio & 51.20 & 82.54 & 0.59706 & 0.62672 & 13.95 & -1.42 \\
\bottomrule\end{tabular}\end{table}

Compared with .35, the higher-floor exposure policy gains 9.24--19.02 Later case points and loses 1.17--1.57 exposure points. The resulting exposure-versus-case contrast remains negative in cases (-26.78 to -16.81 points) and positive in exposure (11.59--12.72 points). All selected Later pointwise risk intervals and 60-condition upper summaries are below $r$, with the largest $U_{60}=.62802431$. However, a .05/60 tail in 4,000 draws corresponds to only about 3.3 draws above the quantile, so this extreme-tail estimate is coarse. It remains a fit-conditional approximate bootstrap summary under shared temporal shocks, not a new coverage theorem. The check supports the measured exchange over the declared floors while exposing the fragility of the particular Descending choice. Union risk and prospective generalization are not evaluated here.
